\section{Recursive Backtracking}

First, we need to rewrite our brute force algorithm to be recursive, as that's
the foundation for the entire algorithm (duh). To do this, we will introduce a
new method \textit{solve()} which will be our recursive solving method. To
also make things simpler, we will extract the equation evaluation into its own
separate method as well.

This means our SAT solver becomes like so:

\begin{lstlisting}
public interface:
    SATSolver(string input)
    bool solveCNF()

private interface:
    void parse(string equation)
    bool solve(bool[] assignment)
    bool evaluate(bool[] assignment)
    int numVars
    int numClauses
    int[][] clauses
    optional bool[] lits
\end{lstlisting}

The flow essentially becomes like this: \textit{solveCNF()} calls
\textit{solve()}, which calls itself over and over and returns whatever solution
it finds back up to \textit{solveCNF()}, who will handle any sort of cleanup.

Why can't we just make \textit{solveCNF()} call itself recursively? Good question.
That's because \textit{solveCNF()} will handle any setup that
\textit{solve()} will need in order to function properly, and it will handle any
cleanup that \textit{solve()} may create.

Let's now discuss how \textit{solve()} will work.
It will be passed in a \textit{bool[]} for the current assignment of literals.
When the size of our assignment is equal to the number of variables, we will
then evaluate the satisfiability of the assignment, and if it's satisfiable we
return true. Otherwise, we return false, backtrack, retry, until we either find a
satisfiable assignment or we've exhausted them all.

The pseudocode is as so, starting with \textit{evaluate()}:

\begin{algorithm}[H]
\caption{Evaluate CNF}\label{alg:6}
\begin{algorithmic}[1]

\Procedure{evaluate}{$assignment[1..n]$}
\For{each $Clause[1..n]$ in $Clauses$}
    \State $satisfied \gets false$

    \For{$i \gets 0$ up to $n$}
        \State $var \gets Clause[i]$
        \State $idx \gets \Call{abs}{var}$

        \If{$var > 0$}
            \State $satisfied \gets assignment[idx] \vee satisfied$
        \Else 
            \State $satisfied \gets \neg assignment[idx] \vee satisfied$
        \EndIf

        \If{$satisfied = true$}
            \State \texttt{Move onto the next clause}
        \EndIf
    \EndFor
    
    \If{$satisfied = false$}
        \State {return $false$}
    \EndIf
\EndFor
\State $vars \gets assignment$
\State return $true$
\EndProcedure
\end{algorithmic}
\end{algorithm}

For this algorithm, we've extracted out the logic on how we evaluate a CNF equation, 
but there are a few changes to how we've done it in our brute force methods. For instance,
we no longer require a global $solved$ variable for the entire equation, but now we only 
need a $satisfied$ variable for just the current clause. Since CNF is satisfied iff all 
clauses are satisfied, if we find one unsatisfied clause, that's enough to stop evaluation 
completely. Otherwise, we can save whatever assignment it was and return true.

Now for our recursive \textit{solve()} method:

\begin{algorithm}[H]
\caption{Recursive Backtracking}\label{alg:7}
\begin{algorithmic}[1]

\Procedure{solve}{$assignment[1..n]$}

\If{$n = numVars$}
    \State return $\Call{evaluate}{assignment}$
\EndIf

\State $\Call{Append}{assignment, true}$

\If{$\Call{solve}{assignment}$}
    \State return $true$
\EndIf

\State $assignment[n] \gets false$

\If{$\Call{solve}{assignment}$}
    \State return $true$
\EndIf

\State $\Call{truncate}{assignment, 1}$
\State return $false$

\EndProcedure
\end{algorithmic}
\end{algorithm}

This might be slightly confusing at first, so let's go over it.

We first check our base case. If our assignment has all of its variables, we
evaluate it and return the result whether it be true or false. And you'll notice
that whenever we return true from a call of \textit{solve()}, we immediately
return true afterwards, so therefore, if we find a satisfiable assignment, we
will go up our call stack back to \textit{solveCNF()}, so no need to worry that
it will terminate if we find a satisfiable assignment.

However, if our assignment still needs some variables, we append true to our
assignment, then we do a recursive call on \textit{solve()}. If that recursive
call leads to a satisfiable assignment, we return true, otherwise, we swap out
our last value for false and try \textit{solve()} again. And if that ends up
working out, we return true, but if it doesn't then we remove the last value
and return false.

So how do we know that it will guarantee termination for an unsatisfiable problem?
Notice how when assigning the current variable true or false and it doesn't
work, we just remove the last one and go up the call stack by one? Well
suppose we're at the top of the call stack (we've assigned the first variable)
and the first variable is assigned false, and that didn't lead to a satisfiable
assignment. Then we remove that first variable from the assignment, and then
we return false. So our assignment array is empty, and we've returned to
\textit{solveCNF()}. Thus, \textit{solve()} will terminate eventually.

And how do we know that we'll try every single assignment?
Well for the very last variable, it's easy to see, if it being true doesn't satisfy
the equation, we assign it false, and if that doesn't work, we remove it and go
up the call stack. Then for the previous variable, if it was true, we assign it
false and go back to the last variable where it tries true and/or false again.
Then that doesn't work, we go up the call stack twice, and that continues until
we've found a satisfiable assignment or we've exhausted them all.\footnote{
This is a very informal proof, but what's most important is that we build up
our understanding of what's going on instead of dumping notation. Perhaps I
will do a formal proof in the future.}

Lastly, our \textit{solveCNF()} method:

\begin{algorithm}[H]
\caption{Main Solving Method}\label{alg:8}
\begin{algorithmic}[1]

\Procedure{solveCNF}{}

\State assignment[1..n]
\State return $\Call{solve}{assignment}$

\EndProcedure
\end{algorithmic}
\end{algorithm}

Because we're dealing with pseudocode, we don't have to do things such as
reserving or resizing the assignment array. And since \textit{evaluate()} takes
care of making sure our \textit{lits} field holds a valid assignment (if found)
and holds nothing if not, we don't need to do anything else in
\textit{solveCNF()}.

You may wonder though, why not pass in \textit{lits} to \textit{solve()} or
have no setup and make \textit{solveCNF()} purely recursive?
Those are very valid questions to ask.

This has to do with the differences between pseudocode and actual
implementation. In the actual implementation, some methods depend on the
state of \textit{lits}, such as whether or not it holds anything, and by passing
it into \textit{solve()}, we force it to hold something, even if there is no
satisfiable assignment. And it's true, that in \textit{solveCNF()} if we had no
setup, we could make it recursive on \textit{lits}, however, I chose to
implement it differently.

When you are reading this, please ask yourself these kinds of questions,
because oftentimes, the way you read my implementation is not the best way,
and even if it is, there is no reason as to why you can't think of different
implementations, try them out yourself, and compare. Perhaps, and I expect
this will happen often, you will find a better implementation than the one
written down here.

\subsection{Recursive Backtracking -- Results}

Running this code on the same 1000 instances (uf20-91), it took
\textbf{1,846,644ms}, which is roughly 30.78 minutes. This is actually a pretty big
improvement from our 35.32 minutes via brute forcing. Although this seems
like a big gain in speed, it should be noted that at this point the algorithm is
still just a fancier brute forcing algorithm. The next time I rewrote this algorithm,
it took \textbf{1,540,613ms}, which is about 25.68 minutes. So once again, that's about a 
~5 minute speedup just from rewriting it. Perhaps we should rewrite every piece of 
software in the world to make it run 5 minutes faster. 

Here's a question for you, why might assigning our variables true instead of
false first lead to an increase in speed? \textbf{Hint: Look at the number of negations 
present in the testing set.}

\begin{FVerbatim}
On all 1000 equations in uf20-91:

Recursive Backtracking:    1,540,613ms (~25.68 minutes).
Brute Force:               1,770,851ms (~29.51 minutes).
Short-Circuited (Basic):   49,874ms (~50 seconds).
Parallel:                  10,589ms (~10.5 seconds).
\end{FVerbatim}
